\section{The thrust mixer, from the boat's geometry}
\label{sec:mixer}

The mixer turns a body command (surge, sway, yaw) into one command per
thruster. It looks like a table of signs, but it comes straight out of where
the thrusters are and which way they push. This section derives it, explains
why it is only $\pm1$ for this boat, what changes when the geometry is not
ideal, and what the thrusters do to a command on its way to the water. The
correction columns are defined in Section~\ref{sec:thrust}.

\subsection{The vessel and its thrusters}

All of Part~I works in the body frame of Section~1 (FLU: $x$ forward, $y$ to
port, yaw counter-clockwise seen from above), with the origin at the
\emph{center of mass}. Figure~\ref{fig:corr-geometry} shows the ideal simulator fixture from
above: two hulls, and four thrusters near the corners of a rectangle, each
turned $45^\circ$ off the hull axis. Describe thruster $i$ by
\begin{itemize}[nosep]
  \item its position $(x_i, y_i)$ in the body frame,
  \item the direction its positive thrust pushes the boat, an angle $\alpha_i$
        measured from $+x$ toward $+y$,
  \item its thrust command $t_i\in[-1,1]$ (negative pushes the other way).
\end{itemize}

\begin{figure}[htbp]
\centering
\begin{tikzpicture}[scale=5.2, >=Stealth]
  % hulls (body x up the page, body y to the left)
  \foreach \s in {1,-1}
    \draw[fill=sD!12, draw=sD] (-\s*0.33-0.05,-0.53) rectangle (-\s*0.33+0.05,0.53);
  \draw[sD, fill=sD!6] (-0.25,-0.25) rectangle (0.25,0.25);
  % body axes at the center of mass
  \draw[->, thick] (0,0) -- (0,0.42) node[above] {$x$ (surge)};
  \draw[->, thick] (0,0) -- (-0.42,0) node[left] {$y$};
  \draw[->] (0.06,0) arc (0:300:0.06);
  \node[font=\scriptsize] at (0.1,0.07) {$r$};
  \fill (0,0) circle (0.012);
  \node[font=\scriptsize, anchor=north west] at (0.03,-0.06) {CM};
  % thrusters: name, page position (-y, x), page direction of +thrust (-sin a, cos a)
  \foreach \n/\px/\py/\dx/\dy/\al in {FL/-0.34/0.27/-0.707/-0.707/135,
                                      FR/0.34/0.27/0.707/-0.707/-135,
                                      RR/0.34/-0.27/0.707/0.707/-45,
                                      RL/-0.34/-0.27/-0.707/0.707/45} {
    \draw[thick, fill=white] (\px,\py) circle (0.03);
    \draw[->, sA, very thick] (\px,\py) -- ++(0.16*\dx,0.16*\dy);
    \node[font=\footnotesize] at (\px+0.09*\dx*-1.4,\py+0.09*\dy*-1.4) {\n};
  }
  % line of action of RL and its lever arm (foot of the perpendicular at (-0.305,-0.305))
  \draw[sB, dashed] (-0.34-0.25*-0.707,-0.27-0.25*0.707) -- (-0.34+0.55*-0.707,-0.27+0.55*0.707);
  \draw[sB, thick] (0,0) -- (-0.305,-0.305);
  \node[sB, font=\scriptsize, anchor=north west] at (-0.2,-0.17) {$\ell = \tfrac{a+b}{\sqrt2}$};
  \draw[sB] (-0.305,-0.305) ++(0.028,0.028) -- ++(-0.028,0.028) -- ++(-0.028,-0.028);
  % dimensions
  \draw[|<->|, sD] (0.45,0) -- (0.45,0.27) node[midway, right, font=\scriptsize] {$a$};
  \draw[|<->|, sD] (0,-0.6) -- (-0.34,-0.6) node[midway, below, font=\scriptsize] {$b$};
  \draw[sD, very thin] (-0.34,-0.3) -- (-0.34,-0.62);
  \node[corr note, sA] at (0.0,-0.75) {blue: direction of positive thrust $(\cos\alpha_i,\sin\alpha_i)$};
\end{tikzpicture}
\caption{The ideal fixture from above, bow up. Body $x$ points forward and $y$ to port, so
yaw $r$ is counter-clockwise. Each thruster pushes along a diagonal. The dashed line
is the line of action of RL; its distance $\ell$ from the center of mass is the
lever arm with which RL turns the boat.}
\label{fig:corr-geometry}
\end{figure}

\subsection{From thrusts to a body force}

Thruster $i$ with thrust $T_i$ (newtons) produces the force
$T_i(\cos\alpha_i, \sin\alpha_i)$. Its turning moment about the center of mass
is the 2-D cross product of lever arm and force,
$\tau_i = x_i F_{y,i} - y_i F_{x,i}$: a force to port ($F_y>0$) applied ahead
of the center ($x_i>0$) turns the bow to port, which is positive yaw. Adding the
four thrusters,
\[
  \underbrace{\begin{bmatrix} F_x\\ F_y\\ \tau \end{bmatrix}}_{\bm{w}}
  = \underbrace{\begin{bmatrix}
      \cos\alpha_1 & \cdots & \cos\alpha_4\\
      \sin\alpha_1 & \cdots & \sin\alpha_4\\
      x_1\sin\alpha_1 - y_1\cos\alpha_1 & \cdots & x_4\sin\alpha_4 - y_4\cos\alpha_4
    \end{bmatrix}}_{B\ \in\ \R^{3\times4}}
  \underbrace{\begin{bmatrix} T_1\\ \vdots\\ T_4 \end{bmatrix}}_{\bm{T}} .
\]
Column $i$ of $B$ is what one newton from thruster $i$ does to the
boat. $B$ is the boat's \emph{allocation} matrix: kinematics, no tuning. The
three numbers $\bm w = (F_x, F_y, \tau)$ are called a \emph{wrench}: a force and
a moment together.

\subsection{Why this boat's mixer is only signs}

The ideal fixture is a symmetric ``X'': thrusters at $(\pm a, \pm b)$, each pushing
along a diagonal. Matching the code's sign convention,
\begin{center}
\begin{tabular}{lcccc}
  \toprule
  & FL & FR & RR & RL\\
  \midrule
  position & $(a, b)$ & $(a, -b)$ & $(-a, -b)$ & $(-a, b)$\\
  $\alpha_i$ & $135^\circ$ & $-135^\circ$ & $-45^\circ$ & $45^\circ$\\
  \bottomrule
\end{tabular}
\end{center}
Every $\cos\alpha_i$ and $\sin\alpha_i$ is $\pm\frac{1}{\sqrt2}$, and every
lever arm works out to $\pm\frac{a+b}{\sqrt2}$, so
\[
  B = \frac{1}{\sqrt2}\begin{bmatrix}
    -1 & -1 & +1 & +1\\
    +1 & -1 & -1 & +1\\
    (a{+}b) & -(a{+}b) & (a{+}b) & -(a{+}b)
  \end{bmatrix}.
\]
Its rows are orthogonal. The configured mixer uses the corresponding sign
pattern, with one normalized effort column per axis:
\[
  M = \begin{bmatrix}
    -1 & +1 & +1\\
    -1 & -1 & -1\\
    +1 & -1 & +1\\
    +1 & +1 & -1
  \end{bmatrix}
  \begin{matrix}\text{FL}\\ \text{FR}\\ \text{RR}\\ \text{RL}\end{matrix}
\]
With identity correction, read a column to see how a single axis is made (Figure~\ref{fig:corr-columns}): pure
surge, $\bm e = (1,0,0)$, gives $\bm t = (-1,-1,+1,+1)$; each of those thrusts
points forward with a sideways part, and the sideways parts cancel in pairs.
Pure yaw drives diagonal pairs against each other, so the forces cancel and only
a moment remains. The orthogonal columns are why $M^\top M = 4I$: asking for
surge never adds sway or yaw. Two more facts about $M$ are used again and again
in Part~I:
\begin{enumerate}[nosep]
  \item \emph{Every column sums to zero}, $\sum_i M_{ij} = 0$: every axis is made
        by two thrusters running forward and two running in reverse.
  \item \emph{The product of any two columns is minus the third}, entry by entry:
        $M_{i1}M_{i2} = -M_{i3}$ and its permutations. Check it on FL:
        $(-1)(+1) = -(+1)$.
\end{enumerate}

\begin{figure}[htbp]
\centering
\begin{tikzpicture}[scale=2.3, >=Stealth]
  \foreach \col/\name/\shift/\mFL/\mFR/\mRR/\mRL in {1/surge $+$/0/-1/-1/1/1,
                                                      2/sway $+$/1.35/1/-1/-1/1,
                                                      3/yaw $+$/2.7/1/-1/1/-1} {
    \begin{scope}[xshift=\shift cm]
      \foreach \s in {1,-1}
        \draw[fill=sD!12, draw=sD] (-\s*0.33-0.05,-0.53) rectangle (-\s*0.33+0.05,0.53);
      \fill (0,0) circle (0.02);
      \draw[->, sA, very thick] (-0.34,0.27) -- ++(\mFL*-0.18,\mFL*-0.18);
      \draw[->, sA, very thick] (0.34,0.27) -- ++(\mFR*0.18,\mFR*-0.18);
      \draw[->, sA, very thick] (0.34,-0.27) -- ++(\mRR*0.18,\mRR*0.18);
      \draw[->, sA, very thick] (-0.34,-0.27) -- ++(\mRL*-0.18,\mRL*0.18);
      \node[font=\small] at (0,-0.7) {\name};
      \node[font=\scriptsize, sD] at (0,-0.85) {$(\mFL,\mFR,\mRR,\mRL)$};
    \end{scope}
  }
  \draw[->, sB, very thick] (0,0) -- (0,0.3);
  \draw[->, sB, very thick] (1.35,0) -- (1.05,0);
  \draw[->, sB, very thick] (2.7+0.15,0) arc (0:270:0.15);
\end{tikzpicture}
\caption{The three columns of $M$ drawn as thrusts (blue) on the boat of
Figure~\ref{fig:corr-geometry}, with the resulting body motion (orange). Sideways
parts cancel in surge, fore-aft parts cancel in sway, and in yaw all forces cancel
and only the moment of the four equal lever arms remains.}
\label{fig:corr-columns}
\end{figure}

\paragraph{Combined axes.} Each axis is corrected separately with its signed
column $\bm c_j^{s_j}$ from Section~\ref{sec:thrust}, where $s_j$ is the sign
of $e_j$ (zero selects the positive column). For thruster $i$, the mixer sums
these single-axis thrusts and inverts the thrust curve:
\[
  q_{ij}=\sum_k M_{ik}c^{s_j}_{kj}e_j,\qquad
  t_i=\varphi_i^{-1}\!\left(\sum_j\varphi_i(q_{ij})\right),\qquad
  \varphi_i(t)=\begin{cases}t^2,&t\ge0,\\-\rho_i t^2,&t<0.\end{cases}
\]
Thus combined commands superpose the corrected single-axis wrenches before
saturation. The configured reverse ratio is $1$ by default.

\paragraph{Desaturation.} When any effort exceeds full scale, all four are
scaled together:
\[
  \bm t'=\frac{\bm t}{\max(1,\max_i|t_i|)}.
\]
The quadratic model scales every thrust by the square of the same factor,
keeping the wrench direction.
\begin{example}
With identity correction and $\rho_i=1$, $\bm e=(1,1,1)$ gives
$\bm t=(1,-\sqrt3,1,1)$ and
$\bm t'=(1/\sqrt3,-1,1/\sqrt3,1/\sqrt3)$. The delivered wrench is one third
of the unsaturated sum.
\end{example}

\subsection{Other angles and rectangles}

\paragraph{Still mirror-symmetric.} Keep the X mirrored but let the thrusters
sit at any angle $\alpha$ from the long axis and at any rectangle $(\pm a,
\pm b)$. The surge row becomes $\pm\cos\alpha$, the sway row $\pm\sin\alpha$,
and every lever arm $\pm(a\sin\alpha + b\cos\alpha)$. The rows stay
orthogonal, so the $\pm1$ mixer still decouples the axes exactly; only the
\emph{strength} of each axis changes: surge capacity scales with $\cos\alpha$,
sway with $\sin\alpha$, yaw with the lever arm. The ideal fixture uses such a
rectangle: about $a \approx 0.27$\,m fore and aft and $b \approx 0.34$\,m to
the sides (hull offset plus mount), a yaw lever arm of
$(a+b)/\sqrt2 \approx 0.43$\,m. Unequal axis strengths need no new mixer:
they are measured as the speed curves of Section~\ref{sec:thrust}, and the
velocity controller divides them out.

\paragraph{Asymmetric.} Symmetry breaks when the angles differ between
thrusters, the thrusters are not mirrored, the thrusters are not equally strong,
or, most often, the center of mass is not at the geometric center (batteries, a
camera mast). Positions are measured from the center of mass, so moving it by
$(x_c, y_c)$ changes every lever arm to
$(x_i - x_c)\sin\alpha_i - (y_i - y_c)\cos\alpha_i$. The yaw row is then no
longer orthogonal to the surge and sway rows, and the $\pm1$ mixer leaks: a pure
sway command produces the moment
\[
  \tau = \sum_i T_i\,\big[(x_i - x_c)\sin\alpha_i - (y_i - y_c)\cos\alpha_i\big]
  \neq 0 \quad\text{for } \bm t = M(0,1,0)^\top,
\]
so a longitudinal center-of-mass shift can make sway turn the boat. The driver
accepts configured mixer rows and measured correction columns (Section~\ref{sec:thrust}).
\subsection{Thrust is not linear in the command}
\label{sec:corr-phi}

The motor effort $t_i$ becomes a pulse width (Section~\ref{sec:esc}). The
mixer models normalized thrust by $\varphi_i$ above; its inverse is
\[
  \varphi_i^{-1}(f)=\sgn(f)\sqrt{\frac{|f|}{\rho_i\ \text{if }f<0\text{, else }1}}.
\]
This model is homogeneous of degree two:
$\varphi_i(\lambda t)=\lambda^2\varphi_i(t)$ for $\lambda>0$.
Physical thrust also depends on the motor model and its configurable pulse deadband.

\paragraph{The normalized wrench.} Feeding the mixer output through the curve
and back through the geometry gives the wrench the boat actually feels. Because
$B$ is $M^\top$ up to a fixed scale per row, it is convenient to drop the
newtons and the lever arm and define the \emph{normalized wrench}
\[
  \Wn(\bm e)=\tfrac14 M^\top\sum_j\varphi(M\bm c_j^{s_j}e_j),
\]
where $\varphi$ acts on each entry, before desaturation. For the ideal
geometry with a common forward scale $T_{\max}=51.5$\,N, the physical wrench is $S\,\Wn$ with
\[
  S = \diag\big(4T_{\max}/\sqrt2,\ 4T_{\max}/\sqrt2,\ 4T_{\max}\ell\big)
  \approx \diag(146\,\text{N},\ 146\,\text{N},\ 89\,\text{N\,m}),
\]
and the $\tfrac14$ makes $\Wn_j = 1$ the wrench of four thrusters at full forward thrust each
contributing to axis $j$.

\begin{example}[one axis]
With identity correction and uniform $\rho$, command $\bm e = s\,a\,\bm e_j$: axis $j$ only, amplitude $a > 0$, sign
$s = \pm1$. Thruster $i$ gets $t_i = s a M_{ij}$, so two thrusters run forward
at $a$ and two in reverse at $a$ (fact 1 above). The forward ones contribute
$a^2 M_{ij}$ each and the reverse ones $-\rho a^2 M_{ij}$, and every term has the
sign of $s$:
\[
  \Wn_j = \tfrac14\big(2a^2 + 2\rho a^2\big)s = s\,\frac{1+\rho}{2}\,a^2 .
\]
At the default $\rho=1$, full surge is $\Wn_x=1$, about $146$\,N
for this ideal scale.
\end{example}

\paragraph{What the column sums buy.} Two consequences of the structure of $M$
are worth stating before any data, because they say which faults can and cannot
couple the axes.
\begin{itemize}
  \item \emph{A uniform reverse weakness couples nothing.} If every thruster has
        the same $\rho$, the cross wrench of a one-axis command is, for $k\ne j$,
        $\Wn_k = \tfrac14\sum_i M_{ik}\varphi(saM_{ij})$. Each $\varphi$ term is
        $a^2$ times $1$ or $-\rho$, depending only on the sign of $sM_{ij}$, so
        the sum splits into $\tfrac{a^2}{4}(\sum_{\text{fwd}}M_{ik} -
        \rho\sum_{\text{rev}} M_{ik})$. Writing the forward set as the $i$ with
        $sM_{ij} = +1$, $\sum_{\text{fwd}} M_{ik} = \tfrac12\sum_i (1 + sM_{ij})M_{ik}
        = \tfrac12(\sum_i M_{ik} + s\sum_i M_{ij}M_{ik}) = 0$ by fact 1 and the
        orthogonality of the columns; the reverse set the same. So the main
        gains are equal for both signs and there is no leak.
  \item \emph{A single weak reverse leaks with a fingerprint.} If only thruster
        $i$ is weak in reverse, say by a relative-strength factor $\eta_i^-=k_i^-/k_{\rm nominal}^- < 1$, its thrust changes by
        $(1-\eta_i^-)\rho a^2$ in exactly the directions where it runs reversed,
        $sM_{ij} = -1$: one sign of each axis, so three of the six directions.
        There the cross wrench is $\tfrac14(1-\eta_i^-)\rho a^2\,M_{ik}$, the row of
        $M$ for that thruster. This is the ideal one-axis fault pattern for that thruster.
\end{itemize}

\file{mhseals\_hardware/thruster\_mixer.py (mix\_thrusters, decouple); crane\_ml MarineSurfaceProfile.cs, ThrusterActuator.cs}
